\documentclass[../../main.tex]{subfiles}
\graphicspath{{../../image/}{../../figure/}}
\begin{document}

\section{一维空间中的平衡律}\label{section:ch07-一维平衡律}

当 $m=1$ 时，一般平衡律方程组\eqref{align:ch03-平衡律一般形式}
退化为
\begin{align}
\partial_t H(U(x,t),x,t)+\partial_x F(U(x,t),x,t)=\Pi(U(x,t),x,t).
\label{align:ch07-平衡律1d}
\end{align}

方程组\eqref{align:ch07-平衡律1d}自然产生于管道中气体流动、弹性杆或弹性弦的振动等问题的研究，其中介质本身被模拟为一维的。最简单的例子是齐次守恒律方程组，起始于标量守恒律
\begin{align}
\partial_t u+\partial_x f(u)=0.
\label{align:ch07-标量守恒律}
\end{align}

尽管标量守恒律表面上简单，却能为物理学及其他领域中的复杂过程提供有价值的洞见。一段公路上交通流的简单流体动力学理论即是一个恰当的例子。

位置 $x$、时刻 $t$ 的交通状态由交通密度 $\rho(x,t)$（譬如以每英里车辆数计）和交通速度 $v(x,t)$（以每小时英里计）描述。场 $\rho$ 与 $v$ 由车辆守恒律相联系，该守恒律与直线运动的质量守恒\eqref{align:ch02-2-3-2}
相同：
\begin{align}
\partial_t\rho+\partial_x(\rho v)=0.
\label{align:ch07-交通-守恒}
\end{align}
该方程进而由一条行为假设封闭：驾驶员根据局部密度设定车速，$v=g(\rho)$。为了反映拥堵效应，$g$ 必须随 $\rho$ 递减，例如 $g(\rho)=v_0(1-\rho/\rho_0)$，其中 $v_0$ 是限速，$\rho_0$ 是饱和密度，超过它交通便爬行至停滞。对于这样的 $g(\rho)$，\eqref{align:ch07-交通-守恒}变为
\begin{align}
\partial_t\rho+\partial_x\left[v_0\rho\left(1-\frac{\rho}{\rho_0}\right)\right]=0.
\label{align:ch07-交通-模型}
\end{align}

然而，这个朴素模型确实抓住了拥堵公路上交通流的某些定性特征，并成为更精致模型的跳板，后者见\cref{section:ch07-注记}所引文献。


标量守恒律的另一个重要例子是 \emph{Buckley--Leverett 方程}
\begin{align}
\partial_t u+\partial_x\left[\frac{u^2}{u^2+a(1-u)^2}\right]=0,
\label{align:ch07-BuckleyLeverett}
\end{align}
其中 $a$ 是某个正参数。它为两种不混溶流体（相）（如油与水）通过多孔介质的直线流动提供了简单模型，因而在石油工业的强化采油作业中获得应用。变量 $u$ 取值于 $[0,1]$，度量某一流体相的饱和度（即体积分数），而通量度量该相的分数流量率。于是该方程表达质量守恒。若忽略惯性与毛细作用，分数流量率由 Darcy 定律确定，它依赖于两相的粘度比以及相对渗透率之比。反过来，相对渗透率依赖于饱和度 $u$，因而方程得以封闭。\eqref{align:ch07-BuckleyLeverett}中采用的经验通量函数抓住了分数流量率的显著特征。应当注意，与\eqref{align:ch07-交通-模型}相对照，\eqref{align:ch07-BuckleyLeverett}中通量的二阶导数在 $[0,1]$ 上变号。正如我们稍后将看到的，这使得解的结构大为复杂。

标量守恒律的又一个富有教益的例子产生于色谱。化学物质的稀溶质以速度 $v$ 运动，通过细小颗粒固体床的间隙并被吸附在固体表面上，其浓度 $c$ 满足守恒律
\begin{align}
\partial_x c+\partial_t[v^{-1}g(c)]=0,
\label{align:ch07-色谱}
\end{align}
其中 $x$ 是沿床的空间变量，$t$ 是时间，$g(c)$ 是\emph{柱等温函数（column isotherm function）}。读者应当注意，在此例中空间与时间互换了角色。

标量守恒律\eqref{align:ch07-标量守恒律}的一个有趣例子——其中两个变量都是空间变量——产生于由不可压缩基体（如橡胶）以不可拉伸纤维增强的复合材料理论。在数学建模中，假定每一根纤维从基体的任一质点穿过，因而材料沿该纤维切向的方向是不可拉伸的。在此材料构成的物体的平衡态中，假定诸纤维为平行直线，切于某固定向量 $A\in\mathbb{R}^3$。以此平衡态为参考构形，我们考虑一个位移（bilipschitz 同胚）$X^*=X^*(X)$，其变形梯度为 $F=\partial X^*/\partial X$。每根纤维成为一条 Lipschitz 曲线，具有切向量场 $A^*=FA$。由于纤维不可拉伸，$|A^*|=1$。进而，由于材料不可压缩，$\det F=1$。于是 $A^{*\top}=(\det F)^{-1}A^\top F^\top$，故 $\operatorname{div}A^{*\top}=0$，这由\cref{theorem:ch01-坐标变换}
得到。现假定 $A^\top=(1,0,0)$，且变形是平面的，$F_{13},F_{23},F_{31},F_{32}$ 全为零。于是 $A^{*\top}=(F_{11},F_{21},0)$，其中 $F_{11}$ 与 $F_{21}$ 仅依赖于 $X^*$ 的前两个坐标 $(x,y)$。假定 $F_{11}>0$，置 $F_{21}=u(x,y)$，并注意到 $|A^*|=1$ 蕴含 $F_{11}=\sqrt{1-u^2}$，我们得出 $\operatorname{div}A^{*\top}=0$ 退化为标量守恒律
\begin{align}
\partial_y u+\partial_x\sqrt{1-u^2}=0.
\label{align:ch07-纤维增强}
\end{align}

热弹性是方程组有趣例子的丰富源泉。在 Lagrange 坐标下，管道内热弹性流体（气体）的直线绝热流动由\eqref{align:ch03-热弹性系统}
的一维版本控制，形式为
\begin{align}
\begin{cases}
\partial_t u-\partial_x v=0\\
\partial_t v+\partial_x p(u,s)=0\\
\partial_t\left[\varepsilon(u,s)+\tfrac12 v^2\right]+\partial_x[v p(u,s)]=0,
\end{cases}
\label{align:ch07-热弹性}
\end{align}
其中 $u$ 是\emph{比容}（由\eqref{align:ch02-2-3-3}
知 $u=1/\rho$），$v$ 表示速度，$\varepsilon$ 是内能，$p$ 代表\emph{压力}（$-p$ 是应力）。注意 $\rho>0$ 将 $u$ 限制为正值。

此处的热力学关系\eqref{align:ch03-热弹性本构}
读作
\begin{align}
p(u,s)=-\varepsilon_u(u,s),\qquad \theta(u,s)=\varepsilon_s(u,s).
\label{align:ch07-热力学}
\end{align}

方程组\eqref{align:ch07-热弹性}是双曲的，当
\begin{align}
\varepsilon_s(u,s)>0,\qquad \varepsilon_{uu}(u,s)>0,
\label{align:ch07-热弹性双曲}
\end{align}
即绝对温度 $\theta$ 为正且内能 $\varepsilon$ 关于 $u$ 凸，或等价地 $p_u(u,s)<0$。

\eqref{align:ch03-等熵热弹性系统}
的一维版本
\begin{align}
\begin{cases}
\partial_t u-\partial_x v=0\\
\partial_t v-\partial_x\sigma(u)=0,
\end{cases}
\label{align:ch07-p系统}
\end{align}
其中 $\sigma'(u)>0$ 单调递增，是控制气体直线等熵流动的双曲组，也是弹性杆的等熵纵向振动以及弹性层的等熵剪切运动的双曲组。对于气体情形，$u$ 是比容，\eqref{align:ch07-p系统}在取 $\sigma=-p$ 时通常被称为 \emph{$p$ 方程组}。对于杆的情形，$u$ 是\emph{应变}。比容与应变都被限制为正值，而在剪切运动情形 $u$ 可取任意符号值。在下文，我们将常以\eqref{align:ch07-p系统}作为数学范式，并假定 $\sigma$ 是 $(-\infty,\infty)$ 上的光滑单调递增函数。

在 Euler 坐标下，气体的直线等熵流动由\eqref{align:ch03-等熵Euler}
的一维版本控制，即
\begin{align}
\begin{cases}
\partial_t\rho+\partial_x(\rho v)=0\\
\partial_t(\rho v)+\partial_x[\rho v^2+p(\rho)]=0.
\end{cases}
\label{align:ch07-等熵Euler}
\end{align}
当 $p'(\rho)>0$ 时该方程组是双曲的。特别地，当流体是理想气体\eqref{align:ch02-2-5-31}
时，\eqref{align:ch07-等熵Euler}成为
\begin{align}
\begin{cases}
\partial_t\rho+\partial_x(\rho v)=0\\
\partial_t(\rho v)+\partial_x[\rho v^2+\kappa\rho^\gamma]=0.
\end{cases}
\label{align:ch07-理想气体}
\end{align}
当 $\gamma>1$ 时，在真空态 $\rho=0$ 处双曲性失效。

所谓\emph{无压气体动力学（pressureless gas dynamics）}方程组
\begin{align}
\begin{cases}
\partial_t\rho+\partial_x(\rho v)=0\\
\partial_t(\rho v)+\partial_x(\rho v^2)=0,
\end{cases}
\label{align:ch07-无压气体}
\end{align}
不是双曲的，它控制一簇"粘性"颗粒的流动：碰撞的颗粒融合为单个颗粒，合并其质量并以保持总线性动量的速度运动。\eqref{align:ch07-无压气体}的解发展出质量集中的倾向，或许可用来解释宇宙中大尺度结构的形成。

接下来我们推导控制三维齐次热弹性介质（参考密度 $\rho_0=1$）等熵平面振荡的方程组。按第 II 章的术语与记号，我们考虑特别形式 $\chi=x+\phi(x\cdot\nu,t)$ 的运动，其中 $\nu$ 是指向振荡方向的（常）单位向量。为与本章记号一致，我们将标量变量 $x\cdot\nu$ 记为 $x$，于是 $\partial_x=\sum_{\alpha=1}^3\nu_\alpha\partial_\alpha$。$\nu$ 方向的速度为 $v(x,t)=\partial_t\phi(x,t)$。我们还置 $u(x,t)=\partial_x\phi(x,t)$，此时变形梯度为 $F=I+u\otimes\nu$。垂直于 $\nu$ 的平面上每单位面积的应力向量为 $\varsigma(u)=S(I+u\nu^\top)\nu$，其中 $S(F)$ 是 Piola--Kirchhoff 应力。于是我们最终得到一个由六个守恒律构成的方程组
\begin{align}
\begin{cases}
\partial_t u-\partial_x v=0\\
\partial_t v-\partial_x\sigma(u)=0,
\end{cases}
\label{align:ch07-平面振荡}
\end{align}
它看起来与\eqref{align:ch07-p系统}相同，只不过这里 $u,v,\sigma$ 不再是标量而是三维向量。

内能 $\varepsilon(F)$ 也成为 $u$ 的函数：$\varepsilon(I+u\nu^\top)=e(u)$。于是，由\eqref{align:ch02-2-5-30}
得 $\sigma(u)=\partial e(u)/\partial u$。因而 $\sigma(u)$ 的 Jacobi 矩阵是 $e(u)$ 的 Hessian 矩阵，也就是在 $F=I+u\nu^\top$ 处赋值的声学张量\eqref{align:ch03-声张量}
。当函数 $e(u)$ 凸时，方程组\eqref{align:ch07-平面振荡}是双曲的。

如第 2.5 节所解释（回忆\eqref{align:ch02-2-5-25}
），当介质是各向同性固体时，内能仅通过不变量 $|F|$、$|F^*|$ 与 $\det F$ 依赖于 $F$。此处 $F=I+u\nu^\top$，故 $|F|^2=3+2u\cdot\nu+|u|^2$，$|F^*|^2=(u\cdot\nu)^2+4u\cdot\nu+|u|^2$，且 $\det F=1+u\cdot\nu$。因而内能仅依赖于两个变量：$|u|$ 与 $u\cdot\nu$。若此外材料不可压缩，运动学约束\eqref{align:ch02-2-7-1}
变为 $u\cdot\nu=0$，此时内能仅依赖于 $|u|$，$e(u)=h(|u|)$。应力张量现由\eqref{align:ch02-2-7-2}
给出，其中 $p$ 是静水压力。简短计算后，回忆 $\varsigma=S\nu$，我们推出\eqref{align:ch07-平面振荡}取形式
\begin{align}
\begin{cases}
\partial_t u-\partial_x v=0\\
\partial_t v+\partial_x p\nu-\partial_x\left(\dfrac{h'(|u|)}{|u|}u\right)=0.
\end{cases}
\label{align:ch07-不可压缩固体}
\end{align}

然而，不可压缩条件 $u\cdot\nu=0$ 蕴含 $\partial_x v\cdot\nu=0$；我们取 $v\cdot\nu=0$ 以消去 $\nu$ 方向上平凡的刚体运动。于是 $\eqref{align:ch07-不可压缩固体}_2$ 给出 $\partial_x p=0$，从而\eqref{align:ch07-不可压缩固体}退化为
\begin{align}
\begin{cases}
\partial_t u-\partial_x v=0\\
\partial_t v-\partial_x\left(\dfrac{h'(|u|)}{|u|}u\right)=0.
\end{cases}
\label{align:ch07-各向同性不可压缩}
\end{align}

\eqref{align:ch07-各向同性不可压缩}的通量函数中所编码的特殊对称性诱导出丰富的几何结构，这对几何学家是一份礼物，而对不得不处理特殊解析困难的分析学家则是一份负担，其滋味稍后将显现。下一个例子表明同样的对称结构也出现在其他场合。我们推导控制柔性、可伸长弹性弦振荡的方程组。弦的参考构形沿 $x$ 轴，假定为（线性）密度为一的自然状态。运动 $\chi=\chi(x,t)$ 通过\emph{速度} $v=\partial_t\chi$ 与\emph{伸长} $u=\partial_x\chi$ 来监测，二者取值于 $\mathbb{R}^3$ 或 $\mathbb{R}^2$，取决于弦是在三维空间中自由运动还是被约束作平面振荡。弦的张力 $\tau$ 假定仅依赖于 $|u|$，$\tau=\tau(|u|)$，它度量弦的伸长。由于弦不能承受任何压缩，$|u|$ 的自然取值范围是 $[1,\infty)$，且 $\tau$ 假定满足 $\tau(r)>0$，$[\tau(r)/r]'>0$，$r>1$。在 Lagrange 坐标下，$u$ 与 $v$ 之间的相容关系连同动量平衡，给出双曲组
\begin{align}
\begin{cases}
\partial_t u-\partial_x v=0\\
\partial_t v-\partial_x\left(\dfrac{\tau(|u|)}{|u|}u\right)=0,
\end{cases}
\label{align:ch07-弹性弦}
\end{align}
它与\eqref{align:ch07-各向同性不可压缩}相同。

我们的下一个例子是控制浅水中长重力波传播的经典守恒律组。它既可由 Euler 方程的渐近分析导出，也可由质量与动量的总体平衡直接（ab initio）导出。我们这里采用后一种途径。

密度为一的不可压缩、无粘流体在水平底、单位宽度的明渠中等熵地流动。自由面上的大气压力取为零。流动由自由面高度变化所诱导的静水压力梯度驱动。假定渠道沿 $x$ 轴，$y$ 轴竖直向上，底面位于 $x$-$z$ 平面。假定自由面的高度在 $z$ 方向上为常数，因而由仅依赖于 $(x,t)$ 的函数 $h$ 描述。此外，速度向量指向 $x$ 方向，且在渠道的任一横截面上为常数，故其长度同样由 $(x,t)$ 的函数 $v$ 描述。

如第 2.7 节所解释，不可压缩无粘流体的应力张量恰为静水压力 $-pI$。$y$ 与 $z$ 方向的线动量平衡分别给出 $\partial_y p=-g$ 与 $\partial_z p=0$，其中 $g$ 是重力加速度。因而 $p=g[h(x,t)-y]$，$0\leqslant y\leqslant h(x,t)$。关于 $y$ 与 $z$ 积分，我们求得时刻 $t$ 作用于 $x$ 横截面上的总压力为 $P(x,t)=\frac12 gh^2(x,t)$。

我们把渠道中的流动看作连续介质的直线运动，由质量与线动量守恒控制，恰如\eqref{align:ch07-等熵Euler}，只不过此处密度的角色自然地由横截面积 $h$ 扮演，压力的角色由压力 $P$ 扮演。于是我们得到浅水波方程组：
\begin{align}
\begin{cases}
\partial_t h+\partial_x(hv)=0\\
\partial_t(hv)+\partial_x\left(hv^2+\tfrac12 gh^2\right)=0.
\end{cases}
\label{align:ch07-浅水波}
\end{align}
注意\eqref{align:ch07-浅水波}与\eqref{align:ch07-理想气体}相同，其中 $\gamma=2$。

正如我们先前所见，通过多孔介质的两相流动（体积分数为 $u$ 与 $1-u$）由 Buckley--Leverett 方程\eqref{align:ch07-BuckleyLeverett}控制。在 $n$ 相流动、体积分数分别为 $U_1,\cdots,U_n$ 的情形，
\begin{align}
U_1+\cdots+U_n=1,
\label{align:ch07-多相-约束}
\end{align}
质量守恒方程退化为如下形式的方程组：
\begin{align}
\partial_t U_i+\partial_x\left[\frac{c_i U_i^2}{\sum_{j=1}^n c_j U_j^2}\right]=0,\qquad i=1,\cdots,n.
\label{align:ch07-多相流}
\end{align}
注意\eqref{align:ch07-多相-约束}与\eqref{align:ch07-多相流}是相容的。当然，可以借助\eqref{align:ch07-多相-约束}消去一个未知量，从而将方程组\eqref{align:ch07-多相流}的规模减一。例如，当 $n=3$ 且 $c_1=c_2=c_3$ 时，
\begin{align}
\begin{cases}
\partial_t u+\partial_x\left[\dfrac{u^2}{u^2+v^2+(1-u-v)^2}\right]=0\\
\partial_t v+\partial_x\left[\dfrac{v^2}{u^2+v^2+(1-u-v)^2}\right]=0.
\end{cases}
\label{align:ch07-三相流}
\end{align}
这类方程组被石油工业用于采油作业。

具有有趣特征的方程组控制平面电磁波通过特殊各向同性电介质的传播，其中电磁能量仅通过标量 $r=(B\cdot B+D\cdot D)^{1/2}$ 依赖于磁感应 $B$ 与电位移 $D$；即在第 3.3.8 节的记号中，$\eta(B,D)=\psi(r)$，满足 $\psi'(0)=0$，$\psi''(0)>0$，且 $r>0$ 时 $\psi'(r)>0$，$\psi''(r)>0$。沿 $3$ 轴方向传播的波由 Maxwell 方程\eqref{align:ch03-Maxwell}
在 $J=0$ 时的解表示，其中场 $B,D,E,H$ 仅依赖于单个空间变量 $x=x_3$ 与时间 $t$。特别地，\eqref{align:ch03-Maxwell}
蕴含 $B_3=0$ 与 $D_3=0$，因而 $B$ 与 $D$ 应视为 $\mathbb{R}^2$ 中满足下述双曲组的向量
\begin{align}
\begin{cases}
\partial_t B-\partial_x\left[\dfrac{\psi'(r)}{r}AD\right]=0\\
\partial_t D+\partial_x\left[\dfrac{\psi'(r)}{r}AB\right]=0,
\end{cases}
\label{align:ch07-电磁波}
\end{align}
其中 $A$ 是交错 $2\times2$ 矩阵，$A_{11}=A_{22}=0$，$A_{12}=-A_{21}=1$。

回到一般平衡律\eqref{align:ch07-平衡律1d}，我们注意到 $H$ 与/或 $F$ 可能显式依赖于 $x$，以反映介质的非均匀性。例如，通过（缓慢）变化横截面 $a(x)$ 的管道的等熵气体流动由方程组
\begin{align}
\begin{cases}
\partial_t[a(x)\rho]+\partial_x[a(x)\rho v]=0\\
\partial_t[a(x)\rho v]+\partial_x[a(x)\rho v^2+a(x)p(\rho)]=a'(x)p(\rho),
\end{cases}
\label{align:ch07-变截面管道}
\end{align}
控制，在齐次情形 $a=\text{常数}$ 时退化为\eqref{align:ch07-等熵Euler}。另一方面，$H$ 或 $F$ 对 $t$ 的显式依赖（表示介质的"老化"）相当罕见。相反，$\Pi$ 对 $t$ 的依赖并不少见，因为外部强迫一般是时间依赖的。

源项 $\Pi$ 可依赖于状态向量 $U$，以反映弛豫或反应效应。后一情形的一个简单例子由方程组
\begin{align}
\begin{cases}
\partial_t\rho+\partial_x(\rho v)=0\\
\partial_t(\rho v)+\partial_x[\rho v^2+(\gamma-1)c\rho\theta]=0\\
\partial_t[c\rho\theta+\beta\rho z+\tfrac12\rho v^2]+\partial_x[(\gamma c\rho\theta+\beta\rho z+\tfrac12\rho v^2)v]=0\\
\partial_t(\rho z)+\partial_x(\rho z v)=-\delta h(\theta-\theta_i)\rho z,
\end{cases}
\label{align:ch07-燃烧}
\end{align}
它控制管道中可燃理想气体的流动。除密度 $\rho$、速度 $v$ 与温度 $\theta$ 之外，状态向量在此还包含未燃气体的\emph{质量分数} $z$，取值于 $[0,1]$。\eqref{align:ch07-燃烧}中前三个方程表达质量、动量与能量的平衡。如\eqref{align:ch02-2-5-18}
，压力的状态方程为 $p=R\rho\theta=(\gamma-1)c\rho\theta$，其中 $\gamma$ 是绝热指数，$c$ 是比热。另一方面，与\eqref{align:ch02-2-5-19}
不同，内能此处还依赖于 $z$：$\varepsilon=c\theta+\beta z$，其中 $\beta>0$ 是\emph{反应热}（假定放热）。在\eqref{align:ch07-燃烧}中控制反应的第四个方程里，$h$ 是标准 Heaviside 函数（$\zeta<0$ 时 $h(\zeta)=0$，$\zeta\geqslant0$ 时 $h(\zeta)=1$），$\theta_i$ 是\emph{点火温度}，$\delta>0$ 是反应速率。

抓住\eqref{align:ch07-燃烧}主要特征的一个简单模型方程组是
\begin{align}
\begin{cases}
\partial_t(u+\beta z)+\partial_x f(u)=0\\
\partial_t z=-\delta h(u)z,
\end{cases}
\label{align:ch07-燃烧模型}
\end{align}
其中 $u$ 与 $z$ 都是标量变量，$f(u)$ 是严格递增的凸函数。

作为表现弛豫的源项的例子，考虑管道中理想气体二元混合物的等温流动。混合物的两个组分都满足质量与动量的部分平衡律：对 $\alpha=1,2$，
\begin{align}
\begin{cases}
\partial_t\rho_\alpha+\partial_x(\rho_\alpha v_\alpha)=0\\
\partial_t(\rho_\alpha v_\alpha)+\partial_x[\rho_\alpha v_\alpha^2+v_\alpha\rho_\alpha]=\chi_\alpha.
\end{cases}
\label{align:ch07-混合物}
\end{align}
耦合由源项 $\chi_\alpha$ 诱导，它由 $v_1$ 与 $v_2$ 之间的差异引起，把动量从另一组分传递到 $\alpha$ 组分。特别地，$\chi_1+\chi_2=0$。在非等温流动中，耦合由能量平衡律增强。在更精致的混合物建模中，密度梯度与密度一起作为状态变量出现（Fick 定律），此时浓度的二阶空间导数出现在场方程中。这类项诱导扩散，类似于热传导或粘性的效应。然而这里我们将处理简单方程组 $\eqref{align:ch07-混合物}_1$--$\eqref{align:ch07-混合物}_2$，它是双曲的。

为把混合物实现为单一连续介质，方便的做法是把原状态向量 $(\rho_1,\rho_2,v_1,v_2)$ 替换为新状态变量 $(\rho,c,v,m)$，其中 $\rho$ 与 $v$ 是混合物的密度与平均速度，即 $\rho=\rho_1+\rho_2$，$\rho v=\rho_1 v_1+\rho_2 v_2$，$c$ 是第一组分的浓度，即 $c=\rho_1/\rho$，而 $m=(-1)^\alpha\rho_\alpha(v-v_\alpha)$。假定 $\chi_\alpha=\beta\rho_\alpha(v_\alpha-v)=(-1)^\alpha\beta m$，其中 $\beta$ 是正常数。于是可把方程组\eqref{align:ch07-混合物}重写为形式
\begin{align}
\begin{cases}
\partial_t\rho+\partial_x(\rho v)=0\\
\partial_t(\rho c)+\partial_x(\rho c v+m)=0\\
\partial_t(\rho v)+\partial_x\left[\rho v^2+(v_2+(v_1-v_2)c)\rho+\dfrac{m^2}{\rho c(1-c)}\right]=0\\
\partial_t(\rho c v+m)+\partial_x\left[\rho c v^2+2mv+\dfrac{m^2}{\rho c}+v_1 c\rho\right]=-\beta m.
\end{cases}
\label{align:ch07-混合物重写}
\end{align}
事实上，上述方程组中的第二与第四个方程正是 $\eqref{align:ch07-混合物}_1$ 以新状态变量重写，而第一与第三个方程由把 $\eqref{align:ch07-混合物}_1$ 与 $\eqref{align:ch07-混合物}_2$ 的相应方程相加得到。

一维方程组\eqref{align:ch07-平衡律1d}还从多维方程组\eqref{align:ch03-平衡律一般形式}
导出，当存在对称性（平面、柱面、径向等）把空间依赖约化为单一参数时。在此过程中，母多维齐次守恒律组可能产生一维非齐次平衡律组，作为多维几何效应的反映。例如，控制径向等熵气体流动的一维方程组（由齐次 Euler 方程\eqref{align:ch03-等熵Euler}
导出）是非齐次的：
\begin{align}
\begin{cases}
\partial_t\rho+\partial_r(\rho v)+\dfrac{2\rho v}{r}=0\\
\partial_t(\rho v)+\partial_r[\rho v^2+p(\rho)]+\dfrac{2\rho v^2}{r}=0.
\end{cases}
\label{align:ch07-径向流}
\end{align}
特别地，某些多维现象（如波聚焦）可在一维空间框架内研究。

\end{document}
